% chapters/ch2_related_work.tex — 第 2 章 相关工作
%
% ⚠ 待定：第 1 章已有「国内外研究现状」时，是否需要独立的相关工作章，
%    取决于三个研究内容各自相关的文献量。定稿前先决定，见 OUTLINE.md 的「待决问题」。
%    本节按「可独立成章、也可并入第 1 章现状节」的两用粒度写作。

\chapter{相关工作}
\label{ch:related}

\section{三维高斯的致密化判据}
\label{sec:related-densification}

3DGS 用一组各向异性的三维高斯表示场景，训练中靠自适应密度控制（adaptive density
control，ADC）决定地图的放大。原始判据很直接，把每个高斯在可观测视角上的屏幕空间
位置梯度取均值，跨过阈值就克隆或分裂，再按尺度与不透明度剔除退化基元。这个设计把
「该不该长」压缩成一个标量统计。统计口径本身没有区分梯度从哪种像素来。在线增量式
SLAM 里，MonoGS~\cite{MonoGS} 把这套 ADC 原样搬进关键帧循环，地图生长因此继承了它的
全部假设。

静态重建侧的工作很快发现，这个统计口径的每个环节都不是中性的管道。归一化一端，
Pixel-GS~\cite{pixelgs} 指出大高斯在多数视角只覆盖边缘少量像素，按视角简单取平均会把
高梯度摊薄，于是改按覆盖像素数加权。梯度方向一端，AbsGS~\cite{absgs} 把高频细节糊化
归因于梯度碰撞。大高斯在不同视角受到方向相反的位置梯度，均值永远低于阈值，该分裂的
不分裂，改用同向分量作判据。StopThePop~\cite{stopthepop} 从排序入手，分层逐像素重排
消除 popping，顺带发现视图一致性修复后所需高斯数约省一半。分母一端，Rogge
等~\cite{counting_contributions} 证明视锥可见性过滤对遮挡是盲的，被遮挡高斯照样计数，
摊薄信号、饿死遮挡区，于是把统计限制在实际贡献渲染的视角，并用平均贡献量替代原始
视角计数。信号形态一端，ACD-GS~\cite{acdgs} 认为一阶屏幕梯度分不清真实几何不连续与
高频纹理残差，改用二阶深度响应作曲率先验，并对克隆与分裂施加不对称调节。预算与调度
一端，Taming~3DGS~\cite{taming3dgs} 把无限增长改成逐像素光度贡献加权的构造式精确预算，
MCMC 视角~\cite{mcmcgs} 则把克隆与分裂整体改写为保概率的重定位，在固定基元数下完成
采样迁移。这些修正彼此独立，合起来指向同一个事实。致密化统计的口径会系统性地决定
高斯长在哪里，改动任何一个环节都能改变地图的形态。

动态场景给这条判据链补上了时间维度。RobustSplat~\cite{robustsplat} 在离线多视角设定
下指出，致密化过程本身会把瞬态扰动长进地图，动态物体存在越久，为它克隆出的高斯越多。
该方法先用延迟增长稳住静态结构，再做瞬态掩码自举。TAD-GS~\cite{tadgs} 进一步把时间
可见性直接写进判据。短寿命基元因监督稀疏而致密化不足，它按高斯的时间存在跨度调节
阈值。在线增量设定则是另一番情形。SLAM 系统只有因果观测，没有未来帧，MonoGS 一系的
工作大多沿用原版 ADC。更进一步的质疑来自 TopoGS-SLAM~\cite{topogsslam}，它在静态室内
基准上论证精度在中等基元数就趋于饱和，继续残差驱动地增长只是徒增显存，改用固定拓扑加
覆盖率触发的初始化，把建图变成有界维数的问题。离线一侧说动态区域需要时间感知的
致密化，在线静态一侧说致密化可能根本不必要，两个方向挤压着中间地带。在线动态场景里
的致密化决策，也就是观测因果、以跟踪误差为终点、逐事件控制注入量的那类决策，仍然
缺少专门的处理。这是第 5 章研究内容三的切入点。

\section{动态场景中的高斯 SLAM}
\label{sec:related-gs-slam}

把动态物体从轨迹估计里分出去，是动态 SLAM 的传统主线。DynaSLAM~\cite{DynaSLAM} 用语义
分割加多实例跟踪，把运动像素从数据关联中剔除，再对被遮挡的背景做补绘。这一类先掩蔽、
后重建的做法划定了后续十年的问题形态：动态区域不参与位姿估计，地图只描述静态世界。
它的代价也很明确，被剔除的区域不产生任何地图内容，掩码错了就一并抹掉真实结构。

高斯地图改变了这个格局。可微渲染让位姿、深度与地图进入同一个优化循环，SplaTAM~\cite{SplaTAM}
沿渲染误差优化相机，稠密地图随优化自然生长；DGS-SLAM~\cite{DGSSlam} 把高斯表示引入动态
环境，在剔除动态区域的同时保留完整地图。地图从被动产物变成了可优化的变量。新的问题也
随之而来，渲染误差对像素级扰动高度敏感，动态物体造成的光度残差若不处理，会同时污染
位姿和地图。WildGS-SLAM~\cite{wildgsslam} 为像素引入不确定性，让残差大的区域自动降权，
把「信多少」交给数据本身。本文作者此前的工作 RoGS-SLAM~\cite{RoGSSLAM} 沿同一条思路，
用光流一致性构造可靠性权重，在跟踪与建图中对动态观测降权而不做硬剔除，在 BONN 动态
序列上验证了这条可靠性路线的上界。这些工作的共同点是都在处理「观测是否可信」，即可靠
性信号作用在跟踪与优化的入口。

入口之后的环节少有人问。观测被加权或剔除之后，地图依然按静态规则生长，哪些高斯该
克隆、该分裂、该删除，仍然由致密化统计说了算，而这套统计对「像素为什么不可信」没有
概念。第 2.1 节的修正谱系处理的是统计口径本身的偏差，动态在线设定下还多出一层，统计
的输入就被动态残差污染。研究内容一管观测入口，研究内容二管运行后的地图维护~\cite{CompressionSurvey}，
中间的地图生长环节是第 5 章要补的位置。
